\section*{ВВЕДЕНИЕ}
\addcontentsline{toc}{section}{ВВЕДЕНИЕ}

В условиях стремительного роста объёмов цифровой информации особую роль приобретают методы эффективного хранения и передачи мультимедийных данных. Значительную часть таких данных составляют цифровые изображения, которые широко используются в информационных системах, веб-приложениях, мобильных устройствах и средствах передачи данных. При этом исходные изображения обладают большим объёмом, что создаёт дополнительные требования к пропускной способности каналов связи и объёму памяти.

Одним из наиболее распространённых способов уменьшения объёма графических данных является сжатие изображений с потерями, позволяющее существенно сократить размер файла при сохранении приемлемого визуального качества. Наибольшее распространение получил стандарт JPEG, основанный на использовании дискретного косинусного преобразования, квантования и субдискретизации цветовых компонент. Данный стандарт применяется в большинстве цифровых камер, графических редакторов и веб-браузеров, что делает его изучение и практическую реализацию актуальной задачей.

В рамках данной курсовой работы рассматривается процесс программной реализации алгоритма JPEG-сжатия растровых изображений. Особое внимание уделяется поэтапной реализации ключевых компонентов алгоритма, включая преобразование цветового пространства RGB в YCbCr, хроматическую субдискретизацию формата 4:2:0, дискретное косинусное преобразование, квантование коэффициентов и последующее восстановление изображения для оценки визуального качества.

Целью данной курсовой работы является разработка программного средства для сжатия изображений в формате JPEG с возможностью управления параметром качества и визуального анализа результатов сжатия.

Для достижения поставленной цели необходимо решить следующие задачи:
\begin{itemize}
	\item изучить принципы и этапы алгоритма JPEG-сжатия изображений;
	\item реализовать преобразование цветового пространства RGB в YCbCr и обратно;
	\item реализовать дискретное косинусное преобразование и обратное преобразование;
	\item разработать механизм квантования с учётом параметра качества;
	\item реализовать хроматическую субдискретизацию и апсемплинг;
	\item разработать графический интерфейс для взаимодействия с пользователем;
	\item провести тестирование программы и анализ полученных результатов.
\end{itemize}
